% Generated draft: requires mathematical and source review.
\documentclass[11pt]{article}
\usepackage[margin=25mm]{geometry}
\usepackage{fontspec}
\setmainfont{Latin Modern Roman}
\usepackage{amsmath,amssymb,mathtools,unicode-math}
\setmathfont{Latin Modern Math}
\usepackage{enumitem,needspace,xurl,hyperref}
\setlength{\parindent}{0pt}
\setlength{\parskip}{5pt}
\setlength{\emergencystretch}{3em}
\newcommand{\N}{\mathbb{N}}
\newcommand{\Z}{\mathbb{Z}}
\newcommand{\Q}{\mathbb{Q}}
\newcommand{\R}{\mathbb{R}}
\newcommand{\C}{\mathbb{C}}
\newcommand{\F}{\mathbb{F}}
\newcommand{\A}{\mathbb{A}}
\newcommand{\PP}{\mathbb{P}}
\newcommand{\E}{\mathbb{E}}
\newcommand{\eps}{\varepsilon}
\newcommand{\dd}{\,\mathrm{d}}
\begin{document}
\section*{UnsolvedMath ID: 20002885}
\textbf{Source problem code:} catalogue entry
\subsection{Definitions and mathematical statement}
Let \(f:(\C^n,0)\to(\C^n,0)\) be a finite holomorphic map germ, meaning that the induced homomorphism of convergent power-series local rings makes the source ring a finite module over the target ring. Let \((V,0)\) be a reduced complex analytic variety germ. Give its set-theoretic inverse image the reduced structure,
\[
W=\bigl(f^{-1}(V)\bigr)_{\rm red}.
\]
The question proposed by Linda Rothschild is whether smoothness of \(W\) at zero forces smoothness of \(V\) at zero. Smoothness here means that the complex analytic germ is a complex submanifold, or equivalently that its local analytic ring is regular. It is essential that the hypothesis concerns the reduced inverse image, rather than requiring the entire scheme-theoretic pullback to be smooth.

The source records an affirmative answer when \(\dim V=1\), and gives the following positive-characteristic warning. Over an algebraically closed field \(k\) of characteristic two, take
\[
f(z_1,z_2,z_3)=(z_1^2,z_2^2,z_3+z_1z_2),\qquad
V=\{w_3^2=w_1w_2\}\subset k^3.
\]
The pulled-back equation is
\[
(z_3+z_1z_2)^2-z_1^2z_2^2=z_3(z_3+2z_1z_2).
\]
In characteristic two its reduced zero set is the smooth hyperplane \(z_3=0\), although \(V\) is singular. The scheme pullback has equation \(z_3^2=0\) and is nonreduced. Over \(\C\) the two distinct hypersurfaces in the factorization do not give a smooth inverse image at the origin, so this is not a complex-analytic counterexample.

The supplied research aid reports affirmative curve, hypersurface, and geometric-complete-intersection cases, leaving an unrestricted non-complete-intersection remainder unresolved by the sources it says it verified. For irreducible reduced \(V\) and irreducible reduced \(W\), it reports a degree identity. Let \(d=\deg f\), let \(\delta=\deg(f|_W)\) be the generic degree of the finite map onto \(V\), and let \(\nu\) be the generic length of the scheme-theoretic inverse image along its reduced support \(W\), measured over its residue field. Then
\[
d=\nu\delta.
\]
If \(W\) is a smooth hypersurface with local equation \(t=0\), the aid reports that a reduced defining equation of \(V\) pulls back to a unit times \(t^\nu\), and
\[
\operatorname{ord}_W(\det Df)\ge\nu-1.
\]
The aid further reports a necessary-condition sieve for any remaining counterexample: the target would have to be normal, prefactorial, singular and non-lci, in ambient dimension at least four; its reduced inverse image would lie in the ramification divisor; the map would be nontriangular, of composite degree, with non-birational restriction to \(W\). Here non-lci means not a local complete intersection. The precise meanings and hypotheses for ``geometric complete intersection,'' ``prefactorial,'' and triangularity, and the underlying literature references, are not fully supplied, so this sieve is retained as a reported reduction rather than an independently established classification.

The source motivation is a related but distinct question: for a real-analytic submanifold germ \(M\subset\C^n\), when is \(f(M)\) a manifold? It refers to [ER] without identifying that reference.
\subsection{Short English statement}
If the reduced inverse image of a complex variety under a finite holomorphic self-map is smooth, must the variety be smooth? The source distinguishes this complex-analytic problem from a characteristic-two counterexample, and the supplied aid reports special affirmative cases and degree constraints on any remaining counterexample.
\subsection{Sources}
\begin{itemize}
\item Linda Rothschild, \emph{Finite mappings}, Problem 3 in the AIM workshop \emph{Complexity of mappings in CR geometry}, \url{https://aimath.org/WWN/complexitycrmap/complexitycrmap.pdf}.
\item The source cites [ER] for the real-analytic image motivation, without a full bibliographic entry. The research aid refers to verified primary partial results but does not identify their authors, titles or complete hypotheses in this record.
\item Local UnsolvedMath registry, permanent ID 20002885, exact code \texttt{catalogue entry}; canonical location \texttt{aim-several-complex-variables-notes.json}, notes[42], dated 14 August 2026; statement labeled exact.
\item Research attempt 1: partial result, problem partially solved at the run, rigor label proved, novelty confidence low; registry status partially solved. The degree identity, ramification constraint and sieve are supplied machine-generated research aids, not peer reviewed, requiring independent verification before citation.
\end{itemize}
\end{document}
